\setcounter{chapter}{1}
\chapter{Nhiệt phản ứng}

\section{Tiến hành thí nghiệm}

\subsection{Thí nghiệm 1: Xác định nhiệt dung của nhiệt lượng kế}
    \subsubsection{Công thức tính \(\mathbf{m_0c_0}\):}
        \begin{equation}
            m_0 \cdot c_0=m\cdot c \cdot\frac{(T_3-T_1)-(T_2-T_3)}{T_2-T_3}
        \end{equation}
        Trong đó:
        \begin{itemize}
            \item \(\mathbf{m} = \rho \cdot V = 1 \cdot 50 = 50 (\si{\gram}) \) là khối lượng của 50 \si{\milli\liter} nước.
            \item \(\mathbf{c} = 1 \left(\frac{\si{\calorie}}{\si{\gram} \cdot \si{\kelvin}} \right) \) là nhiệt dung riêng của nước.
        \end{itemize}
    \subsubsection{Kết quả thí nghiệm:}
        \begin{center}
            \begin{tabular}{|C{0.25\textwidth}|C{0.6\textwidth}|}
                \hline
                Nhiệt độ & Giá trị \\ \hline
                \(T_1\) (\si{\celsius}) & \textit{30} \\ \hline
                \(T_2\) (\si{\celsius}) & \textit{58} \\ \hline
                \(T_3\) (\si{\celsius}) & 45 \\ \hline
                \(m_0c_0 (\si{\calorie} / \si{\kelvin})\) & 7,692 \\ \hline
            \end{tabular}
        \end{center}
    \subsubsection{Kết quả tính \(\mathbf{m_0c_0}\):}
        Quy đổi về đơn vị chuẩn Kelvin theo công thức \(T(\si{\kelvin}) = t(\si{\celsius}) + 273,15\):
        \begin{equation}
            \begin{aligned}
                T_1 &= 30 + 273,15 = 303,15 (\si{\kelvin}),\\
                T_2 &= 58 + 273,15 = 331,15 (\si{\kelvin}),\\
                T_3 &= 45 + 273,15 = 318,15 (\si{\kelvin})
            \end{aligned}
        \end{equation}
        Vì biểu thức chỉ chứa hiệu nhiệt độ nên độ chênh lệch theo \si{\kelvin} và theo \si{\celsius} có cùng giá trị số.
        \begin{equation}
            \begin{aligned}
                m_0 \cdot c_0 &= m\cdot c \cdot\frac{(T_3-T_1)-(T_2-T_3)}{T_2-T_3} \\
                &= 50 \cdot 1 \cdot \frac{(318,15-303,15)-(331,15-318,15)}{331,15-318,15} \\
                &= 7,692 (\si{\calorie} / \si{\kelvin})
            \end{aligned}
        \end{equation}

\subsection{Thí nghiệm 2: Hiệu ứng nhiệt của phản ứng trung hòa}
    \subsubsection{Phương trình phản ứng:}
        \[ \ce{NaOH + HCl -> NaCl + H2O} \]
    \subsubsection{Kết quả thí nghiệm:}
        \begin{center}
            \begin{tabular}{|C{0.25\textwidth}|C{0.6\textwidth}|}
                \hline
                Nhiệt độ & Giá trị \\ \hline
                \(T_1\) (\si{\celsius}) & 30 \\ \hline
                \(T_2\) (\si{\celsius}) & 30 \\ \hline
                \(T_3\) (\si{\celsius}) & 35 \\ \hline
                \(Q\)(\si{\calorie}) & 293,46 \\ \hline
                \(\Delta H\)(\si{\calorie / \mol}) & -11738,4 \\ \hline
            \end{tabular}
        \end{center}
    \subsubsection{Tính toán \(\mathbf{Q}\) và \(\mathbf{\Delta H}\):}
        Quy đổi về đơn vị chuẩn Kelvin theo công thức \(T(\si{\kelvin}) = t(\si{\celsius}) + 273,15\):
        \begin{equation}
            \begin{aligned}
                T_1 &= 30 + 273,15 = 303,15 (\si{\kelvin}),\\
                T_2 &= 30 + 273,15 = 303,15 (\si{\kelvin}),\\
                T_3 &= 35 + 273,15 = 308,15 (\si{\kelvin})
            \end{aligned}
        \end{equation}
        Vì công thức chỉ phụ thuộc vào độ biến thiên nhiệt độ nên \(\Delta T\) theo \si{\kelvin} và theo \si{\celsius} có cùng giá trị số.
        \begin{equation}
            \begin{aligned}
                Q &= (m_0c_0 + m_{\ce{NaCl}}c_{\ce{NaCl}}) \cdot \left(T_3 - \frac{T_1+T_2}{2}\right) \\
                &= \left(7{,}692 + 51 \cdot 1\right) \cdot \left(308,15 - \frac{303,15+303,15}{2}\right) \\
                &= 293,46 (\si{\calorie})
            \end{aligned}
        \end{equation}
        Trong đó:
        \begin{itemize}
            \item \(m_0c_0 = 7,692 (\si{\calorie} / \si{\kelvin}) \) là nhiệt dung của nhiệt lượng kế
            \item \( m_{\ce{NaCl}} = \rho_{\ce{NaCl}} \cdot V_{\ce{NaCl}} = 1,02 \cdot 50 = 51 (\si{\gram}) \) là khối lượng dung dịch muối
            \item \(c_{\ce{NaCl}}=1 \left(\frac{\si{\calorie}}{\si{\gram} \cdot \si{\kelvin}} \right)\) là nhiệt dung riêng của dung dịch muối
        \end{itemize}
        \begin{equation}
            n_{\ce{NaCl}} = C_{M} \cdot V_{\ce{NaCl}} = 0,5 \cdot 0,05 = 0,025 (\si{\mole})
        \end{equation}
        \begin{equation}
            \Delta H = - \frac{Q}{n_{\ce{NaCl}}} = - \frac{293,46}{0,025} = -11738,4 (\si{\calorie} / \si{\mole})
        \end{equation}
    \subsubsection{Kết luận:}
        Với kết quả \(\Delta H = -11738,4 < 0\) nên phản ứng \textbf{tỏa nhiệt} với nhiệt lượng: \begin{equation}
            Q = 293,46 (\si{\calorie})
        \end{equation}

\subsection{Thí nghiệm 3: Nhiệt hòa tan của \(\ce{CuSO4}\)}
    \subsubsection{Kết quả thí nghiệm:}

    \begin{center}
        \begin{tabular}{|C{0.25\textwidth}|C{0.6\textwidth}|}
            \hline
            Đại lượng & Giá trị \\ \hline
            \(T_1\) (\si{\celsius}) & 30 \\ \hline
            \(T_2\) (\si{\celsius}) & \textit{34,5} \\ \hline
            \(m_{\ce{CuSO4}} \)(\si{\gram}) & 4,00 \\ \hline
            \(Q\)(\si{\calorie}) & 277,614 \\ \hline
            \(\Delta H\)(\si{\calorie / \mol}) & -11104,56 \\ \hline
        \end{tabular}
    \end{center}

    \subsubsection{Tính toán \(\mathbf{Q}\) và \(\mathbf{\Delta H}\):}
        Quy đổi về đơn vị chuẩn Kelvin theo công thức \(T(\si{\kelvin}) = t(\si{\celsius}) + 273,15\):
        \begin{equation}
            \begin{aligned}
                T_1 &= 30 + 273,15 = 303,15 (\si{\kelvin}),\\
                T_2 &= 34,5 + 273,15 = 307,65 (\si{\kelvin})
            \end{aligned}
        \end{equation}
        Vì công thức chỉ phụ thuộc vào độ biến thiên nhiệt độ nên \(\Delta T\) theo \si{\kelvin} và theo \si{\celsius} có cùng giá trị số.
        \begin{equation}
            \begin{aligned}
                Q &= (m_0c_0 + m_{\text{dd }\ce{CuSO4}}c_{\text{dd }\ce{CuSO4}}) \cdot (T_2 - T_1) \\
                &= (7{,}692 + 54 \cdot 1) \cdot (307,65 - 303,15) \\
                &= 277{,}614 (\si{\calorie})
            \end{aligned}
        \end{equation}
        Trong đó:
        \begin{itemize}
            \item \(m_0c_0 = 7,692 (\si{\calorie} / \si{\kelvin})\) là nhiệt dung của nhiệt lượng kế.
            \item \(m_{\text{dd }\ce{CuSO4}} = m_{\ce{H2O}} + m_{\ce{CuSO4}} = 50 + 4 = 54 (\si{\gram})\) là khối lượng dung dịch sau khi hòa tan.
            \item \(c_{\text{dd }\ce{CuSO4}} = 1 \left(\frac{\si{\calorie}}{\si{\gram} \cdot \si{\kelvin}}\right)\) là nhiệt dung riêng gần đúng của dung dịch.
        \end{itemize}
        \begin{equation}
            n_{\ce{CuSO4}} = \frac{m_{\ce{CuSO4}}}{M_{\ce{CuSO4}}} = \frac{4,00}{160} = 0,025 (\si{\mole})
        \end{equation}
        \begin{equation}
            \Delta H = - \frac{Q}{n_{\ce{CuSO4}}} = - \frac{277{,}614}{0,025} = -11104{,}56 (\si{\calorie} / \si{\mole})
        \end{equation}
    \subsubsection{Kết luận:}
        Với kết quả \(\Delta H = -11104,56 < 0\) nên quá trình hòa tan \(\ce{CuSO4}\) là quá trình \textbf{tỏa nhiệt} với:
        \begin{equation}
            Q = 277,614 (\si{\calorie})
        \end{equation}

\section{Trả lời câu hỏi}

\begin{enumerate}
    \item \textbf{\(\Delta H_{\mathrm{th}}\) của phản ứng \(\ce{HCl + NaOH -> NaCl + H2O}\) sẽ được tính theo số mol \ce{HCl} hay \ce{NaOH} khi cho \SI{25}{\milli\liter} dung dịch \ce{HCl} \SI{2}{M} tác dụng với \SI{25}{\milli\liter} dung dịch \ce{NaOH} \SI{1}{M}? Tại sao?}

        Ta có:
        \begin{equation}
            \begin{aligned}
                n_{\ce{NaOH}} &= 1 \cdot 0,025 = 0,025\ (\si{\mole}), \\
                n_{\ce{HCl}} &= 2 \cdot 0,025 = 0,050\ (\si{\mole}).
            \end{aligned}
        \end{equation}

        Phản ứng xảy ra theo tỉ lệ mol \(1:1\):
        \begin{equation}
            \ce{HCl + NaOH -> NaCl + H2O}
        \end{equation}

        Vì \(n_{\ce{NaOH}} = 0,025\ \si{\mole} < n_{\ce{HCl}} = 0,050\ \si{\mole}\), \ce{NaOH} là chất giới hạn, còn \ce{HCl} dư. Do đó, \(\Delta H_{\mathrm{th}}\) của phản ứng phải được tính theo số mol \ce{NaOH} đã phản ứng, tức là theo \(0,025\ \si{\mole}\). Phần \ce{HCl} dư không tham gia phản ứng nên không dùng để tính nhiệt phản ứng.

    \item \textbf{Nếu thay dung dịch \ce{HCl} \SI{1}{M} bằng dung dịch \ce{HNO3} \SI{1}{M} thì kết quả thí nghiệm 2 có thay đổi hay không?}

        Nếu thay \ce{HCl} bằng \ce{HNO3} thì kết quả thí nghiệm 2 hầu như không thay đổi đáng kể. Nguyên nhân là vì cả \ce{HCl} và \ce{HNO3} đều là acid mạnh, phân li gần như hoàn toàn trong nước:
        \begin{equation}
            \ce{HCl -> H+ + Cl-}
        \end{equation}
        \begin{equation}
            \ce{HNO3 -> H+ + NO3-}
        \end{equation}

        Trong cả hai trường hợp, bản chất của phản ứng trung hòa với base mạnh đều là:
        \begin{equation}
            \ce{H+ + OH- -> H2O}
        \end{equation}

        Vì vậy, nhiệt hiệu ứng của phản ứng trung hòa gần như giống nhau. Nếu có sai khác thì chủ yếu là do sai số thực nghiệm nhỏ.

    \item \textbf{Tính \(\Delta H_3\) bằng lý thuyết theo định luật Hess. So sánh với kết quả thí nghiệm. Theo em sai số nào là quan trọng nhất? Còn nguyên nhân nào khác không?}

        Theo định luật Hess:
        \begin{equation}
            \Delta H_3 = \Delta H_1 + \Delta H_2 = -18,7 + 2,8 = -15,9\,\text{kcal/mol}
        \end{equation}

        So với số liệu thực nghiệm trong báo cáo:
        \begin{equation}
            \Delta H_{3,\text{TN}} = -11104,56\ \left(\frac{\si{\calorie}}{\si{\mole}}\right) \approx -11,10\,\text{kcal/mol}
        \end{equation}

        Như vậy, giá trị thực nghiệm lệch so với giá trị lý thuyết.

        Theo em, nguyên nhân gây sai số quan trọng nhất là \ce{CuSO4} khan bị hút ẩm. Chất này rất dễ hấp thụ hơi nước trong không khí, làm cho mẫu không còn hoàn toàn khan; vì vậy khối lượng chất đem thí nghiệm không còn chính xác, kéo theo nhiệt lượng đo được và giá trị \(\Delta H\) bị sai lệch.

        Một số nguyên nhân khác có thể kể đến là:
        \begin{itemize}
            \item Nhiệt lượng kế không cách nhiệt hoàn toàn nên có sự thất thoát nhiệt ra môi trường.
            \item Sai số của nhiệt kế khi đo nhiệt độ.
            \item Sai số khi đong thể tích dung dịch.
            \item Sai số của cân khi cân hóa chất.
            \item Việc giả sử nhiệt dung riêng của dung dịch hoặc của \ce{CuSO4} bằng \(1\ \si{\calorie} / (\si{\gram} \cdot \si{\celsius})\) chỉ là gần đúng.
            \item \ce{CuSO4} có thể chưa tan hoàn toàn, làm cho nhiệt hòa tan thu được nhỏ hơn giá trị thực tế.
        \end{itemize}
\end{enumerate}